\documentclass[10pt,a4paper]{article}
\usepackage{amsmath,amssymb,amsthm,graphicx}
\usepackage{hyperref}
\newcommand{\uses}[1]{}
\theoremstyle{definition}
\newtheorem{definition}{Definition}[section]
\newtheorem{theorem}{Theorem}[section]
\newtheorem{proposition}{Proposition}[section]
\newtheorem{lemma}{Lemma}[section]
\newtheorem{sublemma}{Sublemma}[section]
\newtheorem{corollary}{Corollary}[section]
\newtheorem{remark}{Remark}[section]
\title{Chebotarev density theorem (complete author Blueprint proof chain)}
\author{PNT+ project contributors}
\date{github-AlexKontorovich-PrimeNumberTheoremAnd; Apache-2.0}
\begin{document}
\maketitle
\noindent\textit{Editorial note.} Author Blueprint proof, not a Lean formalization. The native source at this pinned commit contains this proof entirely in blueprint\_comment blocks. Source citations and any author omissions are retained without mathematical repair.\par\medskip
\noindent\textit{Allowed standard prerequisites (editorial).}\begin{itemize}
\item Author citations to Sharifi Sections 7.1-7.2, as retained in the original proof.
\end{itemize}\medskip
\section{The Chebotarev density theorem: the case of cyclotomic extensions}

Throughout this section, $K$ is a number field, $m \geq 1$ is a fixed integer,
$L = K(\mu_m)$, and $G = \mathrm{Gal}(L/K)$.  (In particular $G$ is abelian, and for
$K=\mathbb{Q}$ one recovers $G \cong (\mathbb{Z}/m\mathbb{Z})^\times$.)
Write $\zeta_K$ and $\zeta_L$ for the Dedekind zeta functions of $K$ and $L$, and for
an abelian character $\chi : G \to \mathbb{C}^\times$ write $L(\chi,s)$ for the associated
Artin $L$-function (equivalently, the Hecke $L$-function of the ideal character attached to
$\chi$; cf.\ Notation~7.1.17 and Proposition~7.1.18 of
\url{https://www.math.ucla.edu/~sharifi/algnum.pdf}).

The goal of this section is to prove the Chebotarev density theorem in the cyclotomic case
(Proposition~\ref{Chebotarev-cyclotomic-density} below), following the classical argument via
Artin $L$-functions as in Sharifi, Propositions~7.1.16--7.1.19 and Proposition~7.2.1.  The
abelian and general cases are then reduced to this one in the subsequent sections.
\begin{lemma}[Artin $L$-function as Euler product]\label{Artin-L-euler}
Let $\chi : G \to \mathbb{C}^\times$ be an abelian character.  For $\Re(s) > 1$ one has
\[
L(\chi,s) \;=\; \prod_{\mathfrak{p}} \bigl(1 - \chi(\mathfrak{p})\, N\mathfrak{p}^{-s}\bigr)^{-1}
\;=\; \sum_{\mathfrak{a} \subseteq \mathcal{O}_K} \chi(\mathfrak{a})\, N\mathfrak{a}^{-s},
\]
where the product runs over nonzero prime ideals of $\mathcal{O}_K$, with the convention
$\chi(\mathfrak{p}) = 0$ if $\mathfrak{p}$ ramifies in the fixed field of $\ker\chi$.
\end{lemma}
\begin{proof}
This is the specialisation of the Artin Euler product (Definition~7.1.15 of Sharifi) to
abelian characters; see Proposition~7.1.18 of
\url{https://www.math.ucla.edu/~sharifi/algnum.pdf}.  Absolute convergence for $\Re(s)>1$
follows from comparison with $\zeta_K(s)$.
\end{proof}
\begin{lemma}[Dedekind-factor]\label{Dedekind-factor}
For $\Re(s) > 1$ one has
\[
\zeta_L(s) \;=\; \prod_{\chi : G \to \mathbb{C}^\times} L(\chi,s),
\]
where the product runs over all (necessarily one-dimensional) irreducible characters of $G$.
\end{lemma}
\begin{proof}\uses{Artin-L-euler}
In general, for a Galois extension $L/K$ and characters $\chi$ of irreducible representations of
$\mathrm{Gal}(L/K)$, Proposition~7.1.16 of
\url{https://www.math.ucla.edu/~sharifi/algnum.pdf} gives
$\zeta_L(s) = \prod_\chi L(\chi,s)^{\chi(1)}$ for $\Re(s)>1$.  In the present abelian
(cyclotomic) setting every irreducible character is one-dimensional, so $\chi(1)=1$ and the
exponents disappear.  Alternatively, comparing Euler factors at an unramified prime
$\mathfrak{p}$: the primes of $L$ above $\mathfrak{p}$ contribute the factor
$(1 - N\mathfrak{p}^{-fs})^{-g}$ to $\zeta_L$, while the Artin factors multiply to the same
quantity by the usual identity $\prod_\chi (1 - \chi(\varphi_{\mathfrak{p}}) X) = 1 - X^f$
(with $X = N\mathfrak{p}^{-s}$ and $f$ the residue degree).
\end{proof}
\begin{lemma}[Simple pole]\label{Dedekind-pole}
The Dedekind zeta function $\zeta_L$ admits a meromorphic continuation to a neighbourhood of
the line $\Re(s)=1$ with a single simple pole at $s=1$ (and is otherwise holomorphic and
nonvanishing on that line after removing the pole).  In particular
$\log \zeta_L(s) \sim \log(s-1)^{-1}$ as $s \to 1^+$.
\end{lemma}
\begin{proof}
This is the standard analytic continuation of Dedekind zeta (Theorem~7.1.12 of
\url{https://www.math.ucla.edu/~sharifi/algnum.pdf}): absolute convergence of the ideal
Dirichlet series and Euler product for $\Re(s)>1$, meromorphic continuation across
$\Re(s)>1-[L:\mathbb{Q}]^{-1}$, and a simple pole at $s=1$.  The logarithmic asymptotic is
immediate from the simple pole.
\end{proof}
\begin{lemma}[Continuation of nontrivial Artin $L$-functions]\label{Artin-L-continuation}
Let $\chi : G \to \mathbb{C}^\times$ be a nontrivial character.  Then $L(\chi,s)$ extends to an
analytic function on the half-plane $\Re(s) > 1 - [K:\mathbb{Q}]^{-1}$, and in particular is
holomorphic at $s=1$.
\end{lemma}
\begin{proof}
Following Proposition~7.1.19 of \url{https://www.math.ucla.edu/~sharifi/algnum.pdf}: let $n$
be the order of $\chi$.  Geometry of numbers supplies the estimate that, for each $n$th root of
unity $\zeta$, the number of ideals $\mathfrak{a}$ with $N\mathfrak{a} \leq N$ and
$\chi(\mathfrak{a})=\zeta$ is $CN + O\bigl(N^{1-[K:\mathbb{Q}]^{-1}}\bigr)$ with $C$ independent of
$\zeta$.  Summing against $\zeta$ cancels the main terms, so
$\sum_{N\mathfrak{a}\leq N} \chi(\mathfrak{a}) = O\bigl(N^{1-[K:\mathbb{Q}]^{-1}}\bigr)$.
Lemma~7.1.5 of Sharifi then yields absolute uniform convergence of
$\sum_{\mathfrak{a}} \chi(\mathfrak{a})\, N\mathfrak{a}^{-s}$ on compact subsets of
$\Re(s) > 1 - [K:\mathbb{Q}]^{-1}$.
\end{proof}
\begin{lemma}[Dedekind-nonvanishing]\label{Dedekind-nonvanishing}
For any nontrivial character $\chi$ of $G = \mathrm{Gal}(L/K)$, one has $L(\chi,1) \neq 0$.
Moreover $L(\chi,s)$ does not vanish for $\Re(s)=1$.
\end{lemma}
\begin{proof}\uses{Dedekind-factor, Dedekind-pole, Artin-L-continuation}
For the value at $s=1$: write $\log\zeta_L(t) = \sum_\chi \log L(\chi,t)$ for real $t>1$.
As $t\to 1^+$, the left side is $\sim \log(t-1)^{-1}$ by Lemma~\ref{Dedekind-pole}.  If some
nontrivial $L(\chi,\cdot)$ had a zero of order $m_\chi \geq 1$ at $s=1$, the right side would
behave like $\bigl(1 - \sum_\chi m_\chi\bigr)\log(t-1)^{-1}$ up to a bounded error
(the trivial character contributes the pole of $\zeta_K$, absorbed into the factorisation),
forcing $1-\sum m_\chi \leq 0$, a contradiction.  Hence $L(\chi,1)\neq 0$ for all nontrivial
$\chi$ (Sharifi, Proposition~7.1.19).

For the rest of the line $\Re(s)=1$, $s\neq 1$: adapt the classical nonvanishing argument for
Dirichlet $L$-functions (comparison of $\zeta_L(s)$, $\zeta_L(s)^3\,|L(\chi,s)|^4$, or the
$3+4\operatorname{Re}$ inequality) using the Euler product of Lemma~\ref{Artin-L-euler} and the
absence of poles of nontrivial $L(\chi,\cdot)$ on $\Re(s)=1$ from
Lemma~\ref{Artin-L-continuation}.
\end{proof}
\begin{lemma}[Orthogonality for Frobenius classes]\label{Chebotarev-orthogonality}
Fix $\sigma \in G$.  For every prime ideal $\mathfrak{p}$ of $\mathcal{O}_K$ unramified in $L$,
\[
\sum_{\chi : G \to \mathbb{C}^\times} \chi(\sigma)^{-1}\,\chi(\varphi_{\mathfrak{p}})
\;=\;
\begin{cases}
|G| & \text{if }\varphi_{\mathfrak{p}} = \sigma,\\
0 & \text{otherwise.}
\end{cases}
\]
Equivalently, if $\sigma(\zeta_m) = \zeta_m^a$ with $\gcd(a,m)=1$, the sum equals $|G|$
precisely when $N\mathfrak{p} \equiv a \pmod{m}$.
\end{lemma}
\begin{proof}
Standard character orthogonality on the finite abelian group $G$.  The reformulation in terms
of $N\mathfrak{p} \bmod m$ uses that $\varphi_{\mathfrak{p}}(\zeta_m) = \zeta_m^{N\mathfrak{p}}$
for the cyclotomic character.
\end{proof}
\begin{proposition}[Cyclotomic Chebotarev density]\label{Chebotarev-cyclotomic-density}
For every $\sigma \in G$, the set of prime ideals $\mathfrak{p}$ of $K$ unramified in $L$ with
Frobenius $\varphi_{\mathfrak{p}} = \sigma$ has Dirichlet density $1/|G|$.
\end{proposition}
\begin{proof}\uses{Dedekind-factor, Dedekind-pole, Dedekind-nonvanishing, Artin-L-euler,
Chebotarev-orthogonality}
This is Proposition~7.2.1 of \url{https://www.math.ucla.edu/~sharifi/algnum.pdf}.  For
$\Re(s)>1$ one has $\log L(\chi,s) \sim \sum_{\mathfrak{p}} \chi(\mathfrak{p})\, N\mathfrak{p}^{-s}$.
Summing against $\chi(\sigma)^{-1}$ and applying Lemma~\ref{Chebotarev-orthogonality} yields
\[
\sum_{\chi} \chi(\sigma)^{-1} \log L(\chi,s)
\;\sim\;
|G| \sum_{\varphi_{\mathfrak{p}}=\sigma} N\mathfrak{p}^{-s}.
\]
On the other hand, by Lemmas~\ref{Dedekind-factor} and~\ref{Dedekind-nonvanishing} the left
side is $\sim \log\zeta_K(s) \sim \log(s-1)^{-1}$ as $s\to 1^+$ (only the trivial character
contributes a pole).  Comparing the two asymptotics gives the asserted Dirichlet density.
\end{proof}
\begin{lemma}[PNT for one character]\label{Dedekind-PNT}
For any nontrivial character $\chi$ of $G$,
\[
\sum_{N\mathfrak{p} \leq x} \chi(\mathfrak{p})\, \log N\mathfrak{p} \;=\; o(x)
\qquad(x\to\infty).
\]
(Equivalently, writing $\Lambda_\chi$ for the von Mangoldt-type coefficients of $-\frac{L'}{L}(\chi,s)$,
one has $\sum_{n\leq x} \Lambda_\chi(n) = o(x)$.)
\end{lemma}

\emph{Remark on prime powers.}
Expanding $-\frac{L'}{L}(\chi,s)$ as a Dirichlet series selects coefficients supported on
prime powers $\mathfrak{p}^j$ with weight $\chi(\varphi_{\mathfrak{p}})^j\log N\mathfrak{p}$.
Thus an \emph{exact} identity expressing $-\sum_\chi \chi(\sigma)^{-1}\frac{L'}{L}(\chi,s)$ as a
sum over primes with Frobenius $\sigma$ should use the condition $\varphi_{\mathfrak{p}}^j=\sigma$
rather than $\varphi_{\mathfrak{p}}=\sigma$.  For the density statement in
Proposition~\ref{Chebotarev-cyclotomic-density} this distinction is immaterial: the $j=1$ terms
agree with the Frobenius condition, while the $j\geq 2$ terms converge absolutely for
$\Re(s)>1/2$ and do not affect the residue at $s=1$.
\begin{proof}\uses{Dedekind-nonvanishing, Artin-L-continuation}
By Lemmas~\ref{Artin-L-continuation} and~\ref{Dedekind-nonvanishing}, $L(\chi,s)$ is holomorphic
and nonvanishing on $\Re(s)\geq 1$, so $-\frac{L'}{L}(\chi,s)$ extends continuously to
$\Re(s)\geq 1$.  The claimed prime-sum estimate then follows by the same Wiener--Ikehara /
Ingham contour argument used for Dirichlet $L$-functions in the prime-number theorem in
arithmetic progressions (cf.\ the material already formalised for Dirichlet $L$-functions
earlier in this file).  Passing through $\Lambda_\chi$ (rather than the bare prime sum) makes
the $j\geq 2$ prime-power contributions explicit and shows they are $O(x^{1/2}\log x)$,
hence absorbable in the $o(x)$ error; cf.\ the remark after the statement.
\end{proof}
\section{The Chebotarev density theorem: the case of abelian extensions}

Now let $L/K$ be an arbitrary finite abelian extension with Galois group $G$, and fix
$\sigma \in G$.  The goal is to show that the primes of $K$ with Frobenius $\sigma$ still have
Dirichlet density $1/|G|$, by reducing to the cyclotomic case already treated.
(Cf.\ Theorem~7.2.2, Step~2, of \url{https://www.math.ucla.edu/~sharifi/algnum.pdf}; cyclic
extensions already suffice for the later reduction to the general case.)
\begin{lemma}[Cyclotomic crossing]\label{Chebotarev-abelian-crossing}
Choose an integer $m\geq 1$ not dividing the discriminant of $L/K$, large enough that
$H := \mathrm{Gal}(L(\mu_m)/L) \cong (\mathbb{Z}/m\mathbb{Z})^\times$ via the cyclotomic character
and $\mathrm{Gal}(L(\mu_m)/K) \cong G \times H$.  For $\sigma\in G$ and $\tau\in H$ write
$S_{\sigma,\tau}$ for the set of primes of $K$ unramified in $L(\mu_m)$ with Frobenius
$(\sigma,\tau)\in G\times H$, and $S_\sigma$ for the set of primes of $K$ unramified in $L$
with Frobenius $\sigma$ in $G$.  Then
\[
\delta_{\mathrm{inf}}(S_\sigma) \;=\; \sum_{\tau\in H} \delta_{\mathrm{inf}}(S_{\sigma,\tau}).
\]
\end{lemma}
\begin{proof}
Every prime counted in $S_\sigma$ lifts to primes in the various $S_{\sigma,\tau}$ according to
the Frobenius in the cyclotomic layer; the identity of lower densities is the usual additivity
of Dirichlet densities over a finite partition (up to the finitely many ramified primes).
\end{proof}
\begin{lemma}[Density after cyclotomic crossing]\label{Chebotarev-abelian-density}
In the notation of Lemma~\ref{Chebotarev-abelian-crossing}, if the order of $\tau\in H$ is
divisible by $|G|$, then $\langle(\sigma,\tau)\rangle \cap (G\times\{1\}) = \{1\}$, so
$L(\mu_m)$ is obtained by adjoining $\mu_m$ to the fixed field
$F := L(\mu_m)^{\langle(\sigma,\tau)\rangle}$.  Applying Proposition~\ref{Chebotarev-cyclotomic-density}
to the cyclotomic extension $F(\mu_m)/F$ and transporting densities as in the reduction step
of Theorem~7.2.2 yields $\delta(S_{\sigma,\tau}) = 1/(|G|\,|H|)$.
\end{lemma}
\begin{proof}\uses{Chebotarev-cyclotomic-density, Chebotarev-abelian-crossing}
See Sharifi, Theorem~7.2.2, Step~2 (display after ``Now suppose that $|G|$ divides the order of
$\tau$'').
\end{proof}
\begin{proposition}[Abelian Chebotarev]\label{Chebotarev-abelian}
Let $L/K$ be finite abelian with Galois group $G$.  For every $\sigma\in G$, the Dirichlet
density of primes of $K$ with Frobenius $\sigma$ exists and equals $1/|G|$.
\end{proposition}
\begin{proof}\uses{Chebotarev-abelian-crossing, Chebotarev-abelian-density,
Chebotarev-cyclotomic-density}
Let $H_n \subset H$ be the set of elements whose order is divisible by $n$.  Summing the
densities of Lemma~\ref{Chebotarev-abelian-density} over $\tau\in H_n$ gives
$\delta_{\mathrm{inf}}(S_\sigma) \geq |H_n|/(|G|\,|H|)$.  Choosing $m$ so that $|H_n|/|H|$ is
arbitrarily close to $1$ (possible since there are primes $m\equiv 1\pmod{n^j}$ by the
cyclotomic case already proved, and $|H_n|/|H|$ tends to $1$ as $j\to\infty$), one obtains
$\delta_{\mathrm{inf}}(S_\sigma) \geq 1/|G|$.  Summing over $\sigma\in G$ forces equality
throughout, so each density exists and equals $1/|G|$.
\end{proof}
\section{The Chebotarev density theorem: the general case}

Finally let $L/K$ be an arbitrary finite Galois extension with group $G$, and let $C\subset G$
be a conjugacy class.  The theorem reduces to the abelian (indeed cyclic) case already proved,
by passing to the fixed field of an element of $C$.
\begin{lemma}[Reduction to a cyclic subextension]\label{Chebotarev-reduction}
Fix $\sigma\in C$ and let $E$ be the fixed field of $\langle\sigma\rangle$, so $L/E$ is cyclic of
degree $f = |\langle\sigma\rangle|$.  Write $S$ for the set of primes of $K$ unramified in $L$
whose Frobenius class equals $C$, and $T_\sigma$ for the set of primes $\mathfrak{P}$ of $E$
unramified in $L$ (and lying over $K$) with Frobenius equal to $\sigma$.  Then
\[
\delta(S) \;=\; \frac{f\,|C|}{|G|}\,\delta(T_\sigma),
\]
whenever either density exists.  In particular, the Chebotarev statement for $L/K$ and $C$
follows from the abelian statement for the cyclic extension $L/E$ and the element $\sigma$.
\end{lemma}
\begin{proof}
If $\mathfrak{P}\in T_\sigma$, then $\varphi_{\mathfrak{P}}=\sigma$ fixes $E$, so $\mathfrak{P}$ has
residue degree one over $K$.  There are exactly $|G|/f$ primes of $L$ over $\mathfrak{P}\cap K$,
and their Frobenii are equidistributed among the $|C|$ elements of the conjugacy class $C$;
exactly $|G|/(f\,|C|)$ of them have Frobenius $\sigma$.  Comparing the Dirichlet series
$\sum N\mathfrak{p}^{-s}$ over $S$ with $\sum N\mathfrak{P}^{-s}$ over $T_\sigma$ (and using
$\sum_{\mathfrak{p}} N\mathfrak{p}^{-s} \sim \sum_{\mathfrak{P}} N\mathfrak{P}^{-s}$) yields the
displayed identity.  See Sharifi, Theorem~7.2.2, Step~1.
\end{proof}

\begin{theorem}[Chebotarev density theorem]\label{Chebotarev-general}
Let $L/K$ be a finite Galois extension of number fields with Galois group $G$, and let
$C\subset G$ be a conjugacy class.  The set of prime ideals $\mathfrak{p}$ of $K$ unramified in
$L$ whose Frobenius conjugacy class equals $C$ has Dirichlet density $|C|/|G|$.
\end{theorem}

\begin{proof}\uses{Chebotarev-reduction, Chebotarev-abelian}
Combine Lemma~\ref{Chebotarev-reduction} with Proposition~\ref{Chebotarev-abelian} applied to
the cyclic extension $L/E$: the latter gives $\delta(T_\sigma) = 1/f$, whence
$\delta(S) = |C|/|G|$.
\end{proof}

\end{document}
